% signal_loss_text.tex -- GENERATED, do not edit by hand.
%
% Draft replacement text for the 21 cm signal-loss result. Regenerate with
%     uv run python horizon_position/make_paper_signal_loss_figure.py
% in the mock_analysis repo; every number in blocks 1, 2, 4, 5 and 6 below is
% computed at generation time from the arrays the figures are drawn from,
% so re-running after the 21 cm ensemble changes updates the prose and the
% figures together and they cannot drift apart. Block 3 is the exception and
% says so.
%
% Framing to preserve if these are edited: neither figure is a proposed
% analysis. Both project onto eigenmodes of a *simulated* nominal instrument,
% and block 5 gives the reason that matters -- an unmodelled 1 m displacement
% puts signal-like power in the first mode past the filter. Nothing here may
% be phrased so as to license reading a residual excess as a detection.
%
% Paste the six blocks into rasti_template.tex as marked. Nothing here is
% \input by the paper -- this file is a staging area, not a dependency, and
% nothing in this repo writes to the paper .tex itself.
%
% NOTE: signal_loss.pdf is 7.6in wide and must go in a `figure*'
% (double-column) environment, not the `figure' that foreground_svd_residual
% .pdf used -- at \linewidth in a single rasti column its tick labels render
% at roughly 3pt.


% ===================================================================
% BLOCK 1 -- section "Minimising Covariance with the 21-cm Signal".
% Replaces the sentence beginning "We emphasise that this analysis only
% quantifies the spectral complexity ...".
% ===================================================================

This analysis quantifies the spectral complexity of the beam-weighted
foregrounds; on its own it says nothing about whether the cosmological signal
survives the same filter. We therefore passed an ensemble of 1769 global
21-cm models through the identical projection, so that the retained signal and
the foreground residual are read off the same axes
(Fig.~\ref{fig:singular_values}). We adopt $N=9$ filtered modes as a
reference operating point: it is the smallest $N$ at which the foreground
residual (1.82\,mK) falls below the median retained signal and stays below it
for every larger $N$; at $N=8$ the residual is 7.75\,mK against a
median retained signal of 4.94\,mK. At the reference point the median
model retains 10 per cent of its band RMS, or 3.01\,mK, and
75 per cent of the ensemble retains more than the foreground residual.
Retention correlates with signal amplitude but is not determined by it
(Spearman +0.53 against trough depth): of the 247 models with an
absorption trough of 140--160\,mK, the retained RMS has a median of
6.3\,mK but an interquartile range of 2.6--8.3\,mK.
What separates them at matched depth is trough width, narrower troughs
retaining the larger fraction (Spearman -0.59 between width and
retained fraction), since a broad trough is the more nearly foreground-like
shape.

This operating point is set by the foregrounds alone, which is all that
Fig.~\ref{fig:singular_values} is asked to establish. Instrumental systematics
are folded in later and can only raise it. The 1\,m antenna displacement of
section~\ref{subsec:fwd_modelling} is one such term: its worst-case residual is
3.20\,mK at $N=9$, comparable to the median retained signal there,
and one further mode brings it to 0.33\,mK against 2.17\,mK
retained (Fig.~\ref{fig:horizon_shift}). We report such crossings under a
`stays below' rule -- the smallest $N$ beyond which the floor never rises above
the median again -- because both quantities fall with $N$ and cross more than
once, so a first-crossing count would be optimistic. We stress that $N$ is used
here to characterise spectral overlap, and not as a filter depth we propose to
apply to data and then interpret the residual of: an unmodelled displacement
leaves signal-like power in precisely the first mode such a filter would keep,
as section~\ref{subsec:fwd_modelling} sets out.

Retained RMS understates what such a filter leaves measurable. The projection
discards the components of a signal that lie along the foreground modes and
keeps the orthogonal complement, so the relevant question is not how much
amplitude survives but which models remain distinguishable within the subspace
that does. Restricting to the 1547 ensemble members with absorption
troughs deeper than 50\,mK and scoring over 70--130\,MHz, the median separation
between a pair of models falls from 74\,mK to 6.6\,mK under the
filter, yet 99 per cent of pairs remain separated by more than 1\,mK and
93 per cent by more than 2\,mK. A measurement in the filtered subspace
therefore constrains the signal to a family of models differing by components
that lie along the foreground modes, rather than collapsing the ensemble to an
indistinguishable residual.

These figures are a floor rather than a forecast. The filter used here is
maximally agnostic: its modes are estimated from the antenna temperature
itself, and no knowledge of the beam or the sky enters. EDGES
\citep{2025PASP..137l5002C} and MIST \citep{2024MNRAS.530.4125M} instead divide
out a simulated beam chromaticity correction before fitting, and REACH
\citep{2022JAI....1150001C} marginalises over a parametrised beam within its
forward model. EIGSEP is designed to do likewise, using the beam measurements
of section~\ref{subsec:beam_mapping} and the forward model of
section~\ref{subsec:fwd_modelling}; every mode of chromaticity that is modelled
rather than filtered is one fewer mode removed from the signal. Filtering is
also a hard projection, whereas a joint fit for the signal and the foregrounds
recovers part of what a projection discards, and the rotational degrees of
freedom described below provide further leverage that a per-spectrum filter
cannot use. Finally, the calculation uses a single simulated sky and beam and
contains no noise, so it describes spectral subspace overlap and not
sensitivity; whether a retained amplitude is detectable is set by the thermal
noise and integration time, which we do not model here.
Fig.~\ref{fig:singular_values} should accordingly not be read as a requirement
that EIGSEP calibrate at the millikelvin level to detect a typical signal. It
is the residual left by the most conservative foreground filter available, and
it sets the scale of the improvement that beam knowledge and joint fitting are
required to deliver.


% ===================================================================
% BLOCK 2 -- caption for signal_loss.pdf, replacing the
% foreground_svd_residual.pdf caption. Keep \label{fig:singular_values}:
% the horizon_shift caption already refers to it. Use figure*, not figure.
% ===================================================================

\caption{Signal loss under foreground-mode filtering. An ensemble of
1769 global 21-cm models is passed through the same projection onto the
leading $N$ eigenmodes of the simulated antenna temperature that is applied to
the foregrounds. Panels (a1) and (a2) show a random subsample of the ensemble
before and after filtering $N=9$ modes, each coloured by the band RMS it
retains at that operating point; the colour scale is logarithmic and shared
with panel (b), so a curve's colour and its height in (b) agree. Curve opacity
also rises with retained RMS, so that the sparse high-retention tail is
visible; it is a small minority, 5.5 per cent of models above 10\,mK.
Panel (b) shows retained RMS against the number of modes filtered for every
model, with the foreground residual in black. At $N=9$ the foreground
residual is 1.82\,mK while the median model retains 3.01\,mK, or
10 per cent of its band RMS, and 75 per cent of the ensemble
retains more than the foreground residual. The structure surviving in (a2) is
ringing from projecting onto a truncated smooth basis rather than a residual
absorption trough, so retained RMS should not be read as retained signal shape.
This filter uses no knowledge of the beam or the sky and is therefore the most
conservative case; instrumental systematics are folded in at
Fig.~\ref{fig:horizon_shift}. See the text.}


% ===================================================================
% BLOCK 3 -- section "Minimising Covariance", rotation paragraph.
% Insert after "... we aim to use both to improve constraints."
%
% These numbers are NOT computed by this script. They come from a probe run
% on horizon_chromaticity/output/chromaticity_eigsep.npz (t_sys, 1296 ori x
% 1436 LST x 201 freq) on 2026-08-19, which is not checked in. Re-derive from
% that cube if the ensemble or the beam changes.
%
% Read the "16 versus 10" carefully: both are matched-residual counts at a
% 0.6 mK foreground residual, NOT operating points. The zenith figure happens
% to be 10 because that is where the zenith residual reaches 0.6 mK; the
% operating point in block 1 is N = 9, where the residual is 1.82 mK.
% Do not let a reader conflate the two.
% ===================================================================

Rotation adds spectral diversity to the data, and it is worth asking what that
diversity costs in model complexity. Repeating the eigenmode analysis over the
full drive grid of 1296 pointings (36 elevations $\times$ 36 azimuths) rather
than the zenith pointing alone, the pooled antenna temperature reaches a
foreground residual of 0.6\,mK with 16 modes, against 10 for the zenith
pointing at the same residual. The foregrounds seen across every accessible
orientation therefore occupy a subspace only marginally larger than that of a
single pointing. Realising the benefit of that diversity requires a joint fit
in which the sky and the beam are shared parameters and the rotations are
known; a per-spectrum projection of the kind used in
Fig.~\ref{fig:singular_values} cannot exploit it, because the information lies
in the correlation between pointings rather than within any one spectrum, and
pooling the pointings into the projection basis only enlarges the subspace
being removed. We defer this analysis to future work.


% ===================================================================
% BLOCK 4 -- caption for horizon_shift.pdf, replacing the existing one.
% Keep \label{fig:horizon_shift} and the figure* environment.
%
% The caption being replaced describes a figure that no longer exists: it
% promises a "dashed red line" at 10 mK, which was swapped for the retained
% 21 cm band, and counts crossings as "within six modes", which was a
% first-crossing count. Under the stays-below rule used throughout, the
% per-axis answers are 3 (east), 2 (north) and
% 10 (up).
% ===================================================================

\caption{Change in the simulated antenna temperature,
$\Delta T_{\text{ant}}$, when the suspended antenna is displaced to the east
(left column), north (middle column), and up (right column). The displacement
shifts the local horizon, shown in Fig.~\ref{fig:horizon}(b), thereby changing
the fraction of sky occulted by the canyon walls and hence $T_{\text{ant}}$.
\emph{Top row:} the difference spectra for a $+1$\,m displacement,
one curve per sidereal hour (24 in total), coloured by LST as indicated
by the colour bar. \emph{Bottom row:} the RMS over frequency of those spectra
after filtering the leading $N$ eigenmodes of the unperturbed antenna
temperature -- the same modes as Fig.~\ref{fig:singular_values} -- for
displacements of 0.1, 1 and 10\,m, coloured light to
dark by magnitude, with all 24 LSTs drawn for each. The colour encodings
of the two rows are therefore different: LST above, displacement below, since
after filtering the LSTs collapse into a single bundle. Grey shows the 21 cm
signal retained under the identical projection (5--95 per cent of the model
ensemble of Fig.~\ref{fig:singular_values}, median dashed), which is the
benchmark the systematic is read against; $N=9$ is the operating point
adopted there.

The displacements are large in amplitude -- up to 8.9\,K at 50\,MHz for
the $+1$\,m upward shift -- but foreground-like, with 99.9 per
cent of that power in the two leading foreground modes, so the same low-order
filtering that removes the sky removes almost all of them. What survives is
narrow rather than smooth: after 9 modes, 99 per cent of the
upward displacement's remaining power sits in mode 10 alone, at
3.18\,mK, above that mode's nominal foreground content
(1.71\,mK) and above the median model's 21 cm content there
(1.66\,mK). Reading the bottom row down the magnitude sequence gives
the position requirement directly: for the upward displacement, the smallest
$N$ beyond which the worst-LST residual stays below the median retained signal
is 2 modes at 0.1\,m, 10 at 1\,m --- one
past the operating point --- and 14 at 10\,m. We do not
treat this filter as an analysis that could be run on data and its residual
attributed to cosmology; see the text.}


% ===================================================================
% BLOCK 5 -- section "Forward Modelling", after the existing discussion of
% the horizon_shift results ("... Motion up or down produces the largest
% change in antenna temperature ...").
%
% This is the paragraph that keeps the two figures from being read as a
% proposed pipeline. Both of them project onto eigenmodes of a *simulated*
% nominal instrument; neither is the analysis EIGSEP intends to run, and the
% numbers below are the reason why.
% ===================================================================

Two things follow from Fig.~\ref{fig:horizon_shift}, and they pull in opposite
directions. The reassuring one is that a position error does not introduce a
new class of spectral structure. The induced $\Delta T_{\text{ant}}$ is large
in amplitude, but 99.9 per cent of its power lies in the two leading
eigenmodes of the unperturbed antenna temperature: it is, to that accuracy,
more foreground. Low-order filtering of the kind that motivates the EIGSEP
design in section~\ref{subsec:covariance} therefore removes almost all of it,
and the eastward and northward displacements fall below the median retained
21 cm signal after 3 and 2 modes respectively.

The cautionary one concerns what is left. Filtering 9 modes -- the
operating point Fig.~\ref{fig:singular_values} adopts from the foregrounds
alone -- leaves the upward displacement with 3.20\,mK, of which
99 per cent is concentrated in a single mode, mode 10. That
one mode carries more power from a 1\,m displacement (3.18\,mK) than
it does from the nominal foregrounds (1.71\,mK) or from the median
21 cm model (1.66\,mK); 73 per cent of the ensemble has
less signal in that mode than a 1\,m shift would put there, and the leftover
resembles the retained signal in shape closely enough to be confused with it
(cosine similarity up to 0.99, median 0.63).

The vertical response scales predictably, which turns this into a
specification. Over the two decades simulated, 0.1 to 10\,m, the
residual left at $N=9$ by an upward displacement is proportional to that
displacement to within 2 per cent per decade and symmetric in its sign,
because a vertical shift lowers the horizon by a near-uniform offset
(Fig.~\ref{fig:horizon}b); this is the parallel spacing of the three magnitude
bundles in Fig.~\ref{fig:horizon_shift}. Holding the injected power to a tenth
of the median
retained signal therefore requires the antenna's vertical position to be known
to roughly 0.1\,m, and we adopt that as the position-monitoring
requirement. The horizontal responses are neither linear nor symmetric --- a
horizontal shift changes the horizon by an amount that depends on where the
cliff edges fall in azimuth, which is why those curves in
Fig.~\ref{fig:horizon}(b) are spiky, and why the east and north bundles in
Fig.~\ref{fig:horizon_shift} are not evenly spaced --- but they are also the
smaller terms, and the vertical axis sets the requirement.

The wider point is methodological. An unmodelled displacement of the size we
must anticipate deposits power that is signal-like in both amplitude and
shape, at exactly the mode boundary where filtering stops. A residual left by
a filter of fixed depth is therefore not evidence of a cosmological signal:
excess with respect to a foreground model is only as trustworthy as the
instrument model behind it, a lesson the field has already drawn from the
scrutiny of published detections
\citep{2018Natur.564E..32H, 2019ApJ...874..153B, 2019ApJ...880...26S}.
Figs.~\ref{fig:singular_values} and~\ref{fig:horizon_shift} should accordingly
be read as characterisations of spectral structure -- how many modes the
instrument-weighted foregrounds occupy, and how a position error compares with
the signal within that basis -- and not as a proposed pipeline. EIGSEP's
analysis instead forward-models the instrument (this section), with the
antenna position among the parameters that are fitted and marginalised over
rather than assumed. The differentiable forward model makes that
marginalisation tractable, and the position sensitivities computed here are
what set the priors it needs.


% ===================================================================
% BLOCK 6 -- section "Forward Modelling". Replaces the sentence beginning
% "To quantify the effect of the horizon coupling, we computed the change in
% antenna temperature ... by 1 m to the east, north, and up ...".
%
% Needed because Fig.~\ref{fig:horizon_shift}'s residual row now shows three
% displacement magnitudes, and only 1 m was previously introduced. No new
% simulations: run_sims.py has always produced +/-0.1, 1 and 10 m for each
% axis (see positions.py), and the other magnitudes were simply unplotted.
% ===================================================================

To quantify the effect of the horizon coupling, we computed the change in
antenna temperature, $\Delta T_{\text{ant}}(\nu)$, caused by displacing the
antenna to the east, north, and up, relative to the nominal position used for
the horizon profile of Fig.~\ref{fig:horizon}. Each axis was displaced by
0.1, 1 and 10\,m, spanning the range from a
displacement we expect to be able to monitor to one large enough to be
obviously unacceptable, so that the dependence on displacement can be read
off rather than inferred from a single case.
